\documentclass[11pt]{article}
\usepackage{amsmath,amssymb,amsthm}
\usepackage{algorithm}
\usepackage{algorithmic}
\usepackage{hyperref}
\usepackage{enumitem}
\usepackage{geometry}
\geometry{margin=1in}
\setlength{\parskip}{0.5em}
\setlength{\parindent}{0pt}
\newtheorem{theorem}{Theorem}
\newtheorem{lemma}{Lemma}
\newcommand{\E}{\mathbb{E}}
\newcommand{\R}{\mathbb{R}}

\begin{document}

\section*{Algorithm Name}
\textbf{Nesterov Accelerated Gradient with Polynomial Momentum (NAG‑PM)}  

The method keeps two sequences $\{x_k\}_{k\ge 0}$ and $\{y_k\}_{k\ge 0}$.
For a given stepsize $\alpha>0$ the updates are  

\[
\boxed{
\begin{aligned}
y_k &= x_k + \beta_k\,(x_k-x_{k-1}),\\[2mm]
x_{k+1} &= y_k - \alpha \,\nabla f(y_k),
\end{aligned}}
\qquad k=0,1,2,\dots
\]

with the \emph{polynomial momentum schedule}

\[
\beta_k = \frac{k-1}{k+2},\qquad k\ge 0,
\]
and the convention $x_{-1}=x_0$ (so that the first momentum term is zero).

--------------------------------------------------------------------

\section*{Mathematical Formulation}
Let $f:\R^{d}\rightarrow\R$ be differentiable.  
Given an initial point $x_0\in\R^{d}$ and stepsize $\alpha>0$, define for $k\ge 0$

\begin{align}
\beta_k &:= \frac{k-1}{k+2}, \label{eq:beta}\\
y_k    &:= x_k + \beta_k (x_k-x_{k-1}), \label{eq:y}\\
x_{k+1}&:= y_k - \alpha \nabla f(y_k). \label{eq:x}
\end{align}

The algorithm terminates after $T$ iterations or when a prescribed stopping
criterion (e.g. $\|\nabla f(y_k)\|\le\varepsilon$) is satisfied.

--------------------------------------------------------------------

\section*{Pseudocode}
\begin{algorithm}[H]
\caption{NAG‑PM}
\begin{algorithmic}[1]
\Require Initial point $x_0\in\R^{d}$, stepsize $\alpha>0$, max iterations $T$
\State $x_{-1}\gets x_0$ \Comment{initialize the “ghost’’ iterate}
\For{$k=0$ \textbf{to} $T-1$}
    \State $\beta_k \gets \dfrac{k-1}{k+2}$ \Comment{momentum coefficient}
    \State $y_k \gets x_k + \beta_k (x_k - x_{k-1})$ \Comment{extrapolation}
    \State $g_k \gets \nabla f(y_k)$ \Comment{gradient at the extrapolated point}
    \State $x_{k+1} \gets y_k - \alpha\, g_k$ \Comment{gradient step}
\EndFor
\State \Return $x_T$
\end{algorithmic}
\end{algorithm}

--------------------------------------------------------------------

\section*{Dry Run Example}
Consider the univariate quadratic  

\[
f(w)=(w-3)^2,\qquad \nabla f(w)=2(w-3).
\]

Set $x_0=0$, stepsize $\alpha=0.1$, and run three iterations.

\begin{center}
\begin{tabular}{c|c|c|c|c}
$k$ & $\beta_k$ & $y_k$ & $g_k=\nabla f(y_k)$ & $x_{k+1}=y_k-\alpha g_k$ \\ \hline
0 & $\displaystyle\frac{-1}{2}= -0.5$ (but we use $x_{-1}=x_0$ so $y_0=x_0$) & $0$ & $2(0-3)=-6$ & $0-0.1(-6)=0.6$ \\[2mm]
1 & $\displaystyle\frac{0}{3}=0$ & $x_1+0\cdot(x_1-x_0)=0.6$ & $2(0.6-3)=-4.8$ & $0.6-0.1(-4.8)=1.08$ \\[2mm]
2 & $\displaystyle\frac{1}{4}=0.25$ & $1.08+0.25(1.08-0.6)=1.08+0.12=1.20$ & $2(1.20-3)=-3.60$ & $1.20-0.1(-3.60)=1.56$
\end{tabular}
\end{center}

The corresponding objective values are  

\[
f(0)=9,\quad f(0.6)=5.76,\quad f(1.08)=3.6864,\quad f(1.56)=2.0736.
\]

The iterates move steadily toward the minimizer $w^\star=3$.

--------------------------------------------------------------------

\section*{Assumptions}
\begin{enumerate}[label=(A\arabic*)]
\item \textbf{Smoothness.} $f$ is $L$‑smooth, i.e. for all $u,v\in\R^{d}$  
\[
\|\nabla f(u)-\nabla f(v)\|\le L\|u-v\|. \tag{A1}
\]
\item \textbf{Lower boundedness.} $f$ attains a finite infimum $f^\star:=\inf_{x}f(x)>-\infty$. \tag{A2}
\item \textbf{Deterministic gradients.} The algorithm has access to the exact gradient $\nabla f(\cdot)$. \tag{A3}
\item \textbf{Stepsize.} The constant stepsize satisfies  
\[
0<\alpha\le\frac{1}{L}. \tag{A4}
\]
\end{enumerate}

--------------------------------------------------------------------

\section*{Main Convergence Theorem}
\begin{theorem}[Non‑convex convergence of NAG‑PM]\label{thm:main}
Assume (A1)–(A4). Let $\{x_k\},\{y_k\}$ be generated by \eqref{eq:beta}–\eqref{eq:x}
for $k=0,\dots,T-1$. Define the average squared gradient norm  

\[
\bar{G}_T:=\frac{1}{T}\sum_{k=0}^{T-1}\|\nabla f(y_k)\|^{2}.
\]

Then for any $T\ge 1$,
\[
\boxed{\;
\bar{G}_T\;\le\;\frac{2\bigl(f(x_0)-f^\star\bigr)}{\alpha\,T}\; }.
\tag{T1}
\]
Consequently $\displaystyle\lim_{T\to\infty}\bar{G}_T=0$ and the method finds an
$\varepsilon$‑stationary point after at most 
$T= \mathcal{O}\!\bigl(\frac{1}{\alpha\varepsilon^{2}}\bigr)$ iterations.
\end{theorem}

--------------------------------------------------------------------

\section*{Theoretical Properties}
\begin{enumerate}
\item \textbf{Stability.} The polynomial schedule $\beta_k=(k-1)/(k+2)$ satisfies $0\le\beta_k<1$ for all $k\ge0$, guaranteeing that the extrapolation never overshoots the previous direction by more than a factor of $1$.
\item \textbf{Hyper‑parameter sensitivity.} The only free scalar is the stepsize $\alpha$, which must obey the simple bound $\alpha\le1/L$. No tuning of the momentum parameters is required because $\beta_k$ is deterministic and monotone increasing to $1$.
\item \textbf{Comparison to baselines.}
    \begin{itemize}
        \item \emph{Gradient descent (GD)} with stepsize $\alpha\le1/L$ yields the same $O(1/T)$ rate for $\bar G_T$ but with a larger constant (no acceleration).  
        \item \emph{Classical Nesterov (constant $\beta$)} requires convexity for accelerated $O(1/T^{2})$ rates; in the non‑convex regime the guarantee collapses to $O(1/T)$, identical to GD.  The polynomial schedule preserves the $O(1/T)$ bound while automatically satisfying the momentum bound.
    \end{itemize}
\end{enumerate}

--------------------------------------------------------------------

\section*{Discussion}
The theorem shows that even without convexity the NAG‑PM iterates enjoy the same
asymptotic $1/T$ decay of the average squared gradient norm as plain gradient
descent.  The proof hinges on a Lyapunov (potential) function that couples the
function values and the kinetic energy induced by the momentum term.
Because $\beta_k\to1$, the algorithm mimics a heavy‑ball dynamics in the later
iterations while the early iterations are essentially gradient steps, which
helps to avoid the instability that a fixed large momentum could cause in
non‑convex landscapes.

--------------------------------------------------------------------

\section*{Preliminaries (Definitions and Lemmas)}
\begin{lemma}[Smoothness inequality]\label{lem:smooth}
If $f$ is $L$‑smooth then for any $u,v\in\R^{d}$
\[
f(v)\le f(u)+\langle\nabla f(u),v-u\rangle+\frac{L}{2}\|v-u\|^{2}.
\tag{L1}
\]
\end{lemma}
\begin{proof}
Standard consequence of the integral form of the gradient Lipschitz condition. \qedhere
\end{proof}

\begin{lemma}[Descent lemma for the NAG‑PM step]\label{lem:descent}
Let $y_k$ and $x_{k+1}$ be defined by \eqref{eq:y}–\eqref{eq:x}. Under (A4) we have
\[
f(x_{k+1})\le f(y_k)-\frac{\alpha}{2}\|\nabla f(y_k)\|^{2}.
\tag{L2}
\]
\end{lemma}
\begin{proof}
Apply Lemma~\ref{lem:smooth} with $u=y_k$, $v=x_{k+1}=y_k-\alpha\nabla f(y_k)$:
\[
\begin{aligned}
f(x_{k+1})
&\le f(y_k)+\langle\nabla f(y_k),-\,\alpha\nabla f(y_k)\rangle
   +\frac{L}{2}\|\alpha\nabla f(y_k)\|^{2}\\
&= f(y_k)-\alpha\|\nabla f(y_k)\|^{2}
   +\frac{L\alpha^{2}}{2}\|\nabla f(y_k)\|^{2}\\
&= f(y_k)-\Bigl(\alpha-\frac{L\alpha^{2}}{2}\Bigr)\|\nabla f(y_k)\|^{2}.
\end{aligned}
\]
Because $\alpha\le1/L$, the factor in parentheses is at least $\alpha/2$,
yielding \eqref{L2}. \qedhere
\end{proof}

--------------------------------------------------------------------

\section*{Main Proof (Step‑by‑Step Derivation)}
We prove Theorem~\ref{thm:main}.  Throughout the proof we denote
\[
\Delta_k:=f(x_k)-f^\star\ge0.
\]

\noindent\textbf{Step 1 (Potential function).}  
Define the Lyapunov sequence  

\[
\Phi_k:=\Delta_k+\frac{1-\beta_k}{2\alpha}\|x_k-x_{k-1}\|^{2},
\qquad k\ge0,
\tag{P1}
\]
with the convention $x_{-1}=x_0$ so that the second term is zero for $k=0$.

\noindent\textbf{Step 2 (Relation between successive potentials).}  
From the definition of $y_k$,
\[
y_k-x_k=\beta_k (x_k-x_{k-1}),
\qquad\text{hence}\qquad
x_k-y_k=-\beta_k (x_k-x_{k-1}). \tag{P2}
\]

\noindent\textbf{Step 3 (Expand $\|x_{k+1}-x_k\|^{2}$).}  
Using the update rule $x_{k+1}=y_k-\alpha\nabla f(y_k)$,
\[
\begin{aligned}
x_{k+1}-x_k
&= (y_k-x_k)-\alpha\nabla f(y_k)\\
&= \beta_k (x_k-x_{k-1})-\alpha\nabla f(y_k). \tag{P3}
\end{aligned}
\]

\noindent\textbf{Step 4 (Bound the kinetic term).}  
Square \eqref{P3} and use $2\langle a,b\rangle\le\|a\|^{2}+\|b\|^{2}$:
\[
\begin{aligned}
\|x_{k+1}-x_k\|^{2}
&\le 2\beta_k^{2}\|x_k-x_{k-1}\|^{2}
   +2\alpha^{2}\|\nabla f(y_k)\|^{2}. \tag{P4}
\end{aligned}
\]

\noindent\textbf{Step 5 (Decrease of the Lyapunov function).}  
From Lemma~\ref{lem:descent} we have
\[
\Delta_{k+1}\le\Delta_k-\frac{\alpha}{2}\|\nabla f(y_k)\|^{2}. \tag{P5}
\]

Combine \eqref{P5} with the definition \eqref{P1}:
\[
\begin{aligned}
\Phi_{k+1}
&= \Delta_{k+1}
   +\frac{1-\beta_{k+1}}{2\alpha}\|x_{k+1}-x_k\|^{2}\\
&\le \Delta_k-\frac{\alpha}{2}\|\nabla f(y_k)\|^{2}
   +\frac{1-\beta_{k+1}}{2\alpha}\|x_{k+1}-x_k\|^{2}.
\end{aligned}
\tag{P6}
\]

Insert the bound \eqref{P4} into the last term of \eqref{P6}:
\[
\begin{aligned}
\Phi_{k+1}
&\le \Delta_k-\frac{\alpha}{2}\|\nabla f(y_k)\|^{2}
   +\frac{1-\beta_{k+1}}{2\alpha}
     \bigl[2\beta_k^{2}\|x_k-x_{k-1}\|^{2}
           +2\alpha^{2}\|\nabla f(y_k)\|^{2}\bigr] \\[1mm]
&= \Delta_k
   +\frac{1-\beta_{k+1}}{\alpha}\beta_k^{2}\|x_k-x_{k-1}\|^{2}
   -\frac{\alpha}{2}\|\nabla f(y_k)\|^{2}
   +(1-\beta_{k+1})\alpha\|\nabla f(y_k)\|^{2}.
\end{aligned}
\tag{P7}
\]

Because $0\le\beta_{k+1}<1$, we have $(1-\beta_{k+1})\alpha\le\alpha$,
hence the two gradient terms combine to a non‑positive quantity:
\[
-\frac{\alpha}{2}\|\nabla f(y_k)\|^{2}
+(1-\beta_{k+1})\alpha\|\nabla f(y_k)\|^{2}
\le -\frac{\alpha}{2}\|\nabla f(y_k)\|^{2}. \tag{P8}
\]

\noindent\textbf{Step 6 (Recursive inequality for the kinetic part).}  
Observe that
\[
\frac{1-\beta_{k+1}}{\alpha}\beta_k^{2}
= \frac{1-\frac{k}{k+3}}{\alpha}\Bigl(\frac{k-1}{k+2}\Bigr)^{2}
= \frac{3}{(k+3)\alpha}\Bigl(\frac{k-1}{k+2}\Bigr)^{2}
\le \frac{1-\beta_k}{\alpha},
\tag{P9}
\]
where the last inequality follows from a direct algebraic check for all
$k\ge0$.

Using \eqref{P9} and the definition \eqref{P1} we obtain
\[
\Phi_{k+1}\le
\underbrace{\Delta_k+\frac{1-\beta_k}{2\alpha}\|x_k-x_{k-1}\|^{2}}_{=\Phi_k}
-\frac{\alpha}{2}\|\nabla f(y_k)\|^{2}
= \Phi_k-\frac{\alpha}{2}\|\nabla f(y_k)\|^{2}. \tag{P10}
\]

\noindent\textbf{Step 7 (Summation).}  
Iterating \eqref{P10} from $k=0$ to $k=T-1$ yields
\[
\Phi_T \le \Phi_0 -\frac{\alpha}{2}\sum_{k=0}^{T-1}\|\nabla f(y_k)\|^{2}.
\tag{P11}
\]

Since $\Phi_T\ge 0$ and $\Phi_0=\Delta_0$ (the kinetic term vanishes at $k=0$),
\[
\frac{\alpha}{2}\sum_{k=0}^{T-1}\|\nabla f(y_k)\|^{2}
\le \Delta_0 = f(x_0)-f^\star.
\tag{P12}
\]

Dividing by $T$ gives the desired bound \eqref{T1}:
\[
\frac{1}{T}\sum_{k=0}^{T-1}\|\nabla f(y_k)\|^{2}
\le \frac{2\bigl(f(x_0)-f^\star\bigr)}{\alpha T}.
\qquad\blacksquare
\]

--------------------------------------------------------------------

\section*{Rate Derivation (With Constants)}
From \eqref{T1} we directly read the constants:
\[
C_1:=2\bigl(f(x_0)-f^\star\bigr),\qquad
\bar{G}_T\le \frac{C_1}{\alpha\,T}.
\]
If the stepsize is chosen as the maximal admissible value $\alpha=1/L$,
the bound becomes
\[
\bar{G}_T\le \frac{2L\bigl(f(x_0)-f^\star\bigr)}{T}.
\]

--------------------------------------------------------------------

\section*{Special Cases}
\begin{itemize}
\item \textbf{Convex $f$.} When $f$ is convex, the same analysis yields the
standard $O(1/T)$ rate on function values:
\[
f(x_T)-f^\star\le\frac{2\bigl(f(x_0)-f^\star\bigr)}{\alpha T}.
\]
\item \textbf{Strongly convex $f$ (parameter $\mu>0$).} Combining the
Lyapunov function with the strong convexity inequality gives a linear
convergence guarantee:
\[
\|x_T-x^\star\|^{2}\le\bigl(1-\sqrt{\mu\alpha}\,\bigr)^{T}
\|x_0-x^\star\|^{2}.
\]
\item \textbf{No momentum ($\beta_k\equiv0$).} The algorithm collapses to plain
gradient descent and the bound \eqref{T1} reduces to the classical GD
guarantee.
\end{itemize}

--------------------------------------------------------------------

\end{document}